\subsection*{専門性}

\addcontentsline{toc}{subsection}{1 専門性}

専門性とは、専門分野の知識を持っているだけでなく、その知識を社会的責任のもとで使う態度を含む。ソフトウェア技術者の専門性は、倫理規程と専門職行動規範に従うこと、標準と実務慣行を理解すること、継続的に学ぶこと、利害関係者へ説明できる判断を行うことに表れる。

専門職団体は、知識体系、倫理規程、専門職行動規範、認定や免許の基準を示す。これらの基準は、技術者が共同体の一員として認められるための入口であると同時に、能力や過失を判断する基準にもなり得る。

\subsubsection*{認定・資格証明・資格・免許}

\addcontentsline{toc}{subsubsection}{1.1 認定・資格証明・資格・免許}

認定、資格証明、資格、免許は似ているが、対象と強制力が異なる。初学者は「組織を見るのか、個人を見るのか」「任意なのか、法的に必要なのか」で区別するとよい。

\begin{center}
\small
\par\smallskip\noindent\refstepcounter{table}\label{tab:ch14-03}表 \thetable: 認定・資格証明・資格・免許の整理\par\smallskip
表 \thetable は、認定・資格証明・資格・免許の整理を比較・確認しやすい形で整理したものである。\par\smallskip
\begin{longtable}{p{0.15\linewidth}p{0.19\linewidth}p{0.33\linewidth}p{0.28\linewidth}}
\toprule
区分 & 対象 & 意味 & 実務上の見方 \\
\midrule
\endfirsthead
\toprule
区分 & 対象 & 意味 & 実務上の見方 \\
\midrule
\endhead
認定 & 教育機関・教育課程・組織 & 一定の教育・品質基準を満たしていることを第三者が認める。 & 工学教育プログラムや大学課程の品質を示す。 \\
\midrule
資格証明 & 個人 & 試験や経験要件により、一定水準の知識・能力を持つことを示す。 & 能力を外部に説明する材料になるが、多くの場合は任意である。 \\
\midrule
資格 & 個人 & 特定の能力や条件を満たしている状態を示す。再資格化が不要な場合もある。 & 資格証明に近いが、制度により意味が異なる。 \\
\midrule
免許 & 個人 & 特定の業務を行う権限を政府や法定機関が与える。 & 国や地域によっては、専門技術者として業務を行う条件になる。 \\
\midrule
\bottomrule
\end{longtable}
\normalsize
\end{center}

\par\smallskip
\begin{center}
\centering
\IfFileExists{assets/generated_figures/ch14/fig-ch14-02.png}{\includegraphics[width=0.96\linewidth]{assets/generated_figures/ch14/fig-ch14-02.png}}{\fbox{\parbox[c][48mm][c]{0.9\linewidth}{\centering 画像生成待ち\\認定・資格証明・免許の違いを示す比較図}}}
\par\smallskip
\refstepcounter{figure}\label{fig:ch14-02}図 \thefigure: 認定・資格証明・免許の違いを示す比較図
\end{center}
図 \thefigure は、認定・資格証明・免許の違いを示す比較図を視覚的に整理したものである。
\par\smallskip

\paragraph*{認定}

\addcontentsline{toc}{subsubsection}{1.1.1 認定}

認定とは、組織、教育機関、教育プログラムなどが一定の能力、権限、信頼性、品質を備えていることを認めることである。工学分野では、認定された教育課程を修了することが、技術者が基礎知識を得る主要な方法になる国もある。国際的な枠組みとして、国際エンジニアリング連合や欧州工学教育認定ネットワークが例として挙げられる。

\paragraph*{資格証明と資格}

\addcontentsline{toc}{subsubsection}{1.1.2 資格証明と資格}

資格証明とは、個人が特定の専門領域で一定水準の能力を持つことを確認する仕組みである。通常は試験と経験要件によって取得し、専門職団体などの非政府組織が運営することが多い。ソフトウェア工学では、標準的な実務知識を確認し、キャリア形成を支える目的がある。

資格証明は、専門職標準を満たす力、専門的判断を適用する力、所定の知識、実務経験、実証された方法論の理解を示す。ただし、それは能力を示す一つの手段であり、すべてのソフトウェア技術者が何らかの資格証明を持つわけではない。

\paragraph*{免許}

\addcontentsline{toc}{subsubsection}{1.1.3 免許}

免許とは、ある人に特定の活動を行う権限を与え、その成果に責任を持たせる制度である。通常は政府機関または法定機関が発行する。免許制度の目的は、資格を持たない者による危険な実務から公衆を守ることである。国や地域によっては、免許を持たない者が専門技術者として業務を行ったり、企業が工学サービスを提供したりできない場合がある。

\subsubsection*{倫理規程と専門職行動規範}

\addcontentsline{toc}{subsubsection}{1.2 倫理規程と専門職行動規範}

倫理規程と専門職行動規範は、技術者の実務と意思決定がどのような価値と行動に基づくべきかを示す。ソフトウェア技術者は、顧客や雇用主だけでなく、公衆の健康、安全、福祉を考慮しなければならない。これは、ソフトウェアの失敗が利用者の生活、社会秩序、財産、生命に影響し得るからである。

倫理規程は、社会規範や地域の法律の文脈の中で成立する。複数の義務が衝突する場面、たとえば顧客の利益、利用者の安全、守秘義務、公共の利益がぶつかる場面で、倫理規程は判断の手がかりになる。ACM、IEEE、IFIPなどの専門職団体は、それぞれ倫理規程や専門職行動規範を公開している。

\par\smallskip
\begin{center}
\centering
\IfFileExists{assets/generated_figures/ch14/fig-ch14-03.png}{\includegraphics[width=0.96\linewidth]{assets/generated_figures/ch14/fig-ch14-03.png}}{\fbox{\parbox[c][48mm][c]{0.9\linewidth}{\centering 画像生成待ち\\倫理規程違反の種類を示す図}}}
\par\smallskip
\refstepcounter{figure}\label{fig:ch14-03}図 \thefigure: 倫理規程違反の種類を示す図
\end{center}
図 \thefigure は、倫理規程違反の種類を示す図を視覚的に整理したものである。
\par\smallskip

\subsubsection*{専門職団体の性質と役割}

\addcontentsline{toc}{subsubsection}{1.3 専門職団体の性質と役割}

専門職団体は、実務家と研究者を含む共同体である。専門職団体は、専門分野を定義し、発展させ、必要に応じて規律を与える。ソフトウェア工学では、専門職団体が知識体系、標準、倫理規程、資格、認定、研修、出版、会議を支える。

個人にとって、専門職団体への参加は、知識を更新し、専門分野とのつながりを保ち、専門家ネットワークを広げる手段である。技術の変化が速いソフトウェア分野では、専門職団体を通じた継続学習が特に重要である。

\begin{itemize}[leftmargin=2em]

\item 一般に受け入れられた知識体系を整備し、広める。

\item 免許、資格証明、認定の基盤を提供する。

\item 倫理規程や専門職行動規範を示し、違反への懲戒を行うことがある。

\item 会議、研修、出版、標準化活動を通じて専門分野を発展させる。

\end{itemize}

\subsubsection*{ソフトウェア工学標準の性質と役割}

\addcontentsline{toc}{subsubsection}{1.4 ソフトウェア工学標準の性質と役割}

ソフトウェア工学標準は、開発、保守、支援、品質、要求、設計、テスト、運用など幅広い話題を扱う。標準は、個々の組織や技術者が毎回ゼロから判断しなくてもよいように、過去の経験と合意をまとめたものである。

標準に従うことには二つの価値がある。第一に、成果物や作業の最低限の性質を明確にし、暗黙の前提や過信を減らす。第二に、法的な争いが生じたときに、一般に受け入れられた実務に従っていたことを説明する根拠になり得る。

\par\smallskip
\begin{center}
\centering
\IfFileExists{assets/generated_figures/ch14/fig-ch14-04.png}{\includegraphics[width=0.96\linewidth]{assets/generated_figures/ch14/fig-ch14-04.png}}{\fbox{\parbox[c][48mm][c]{0.9\linewidth}{\centering 画像生成待ち\\知的財産の違いを示す図}}}
\par\smallskip
\refstepcounter{figure}\label{fig:ch14-04}図 \thefigure: 知的財産の違いを示す図
\end{center}
図 \thefigure は、知的財産の違いを示す図を視覚的に整理したものである。
\par\smallskip

\paragraph*{専門職責任}

\addcontentsline{toc}{subsubsection}{1.7.6 専門職責任}

\term{専門職責任}{Professional Responsibility}とは、技術者が専門家として注意義務を果たさなかった場合に問われ得る責任である。ソフトウェア技術者は、顧客や雇用主にサービスを提供するため、専門過誤、過失、能力不足の主張に備えなければならない。

責任に関する主張への防御として、標準と一般に受け入れられた実務に従っていたことを示すことが重要になる。ソフトウェアの安全性や適合性は測定が難しいため、適切な注意を払わなかったと見なされると過失の根拠になり得る。

\paragraph*{法的要求}

\addcontentsline{toc}{subsubsection}{1.7.7 法的要求}

ソフトウェア技術者は、地域、国、国際的な法的枠組みの中で活動する。登録、免許、試験、教育、経験、研修などの要件、契約上の合意、責任に関する非契約上の法律などを理解する必要がある。

\paragraph*{貿易遵守}

\addcontentsline{toc}{subsubsection}{1.7.8 貿易遵守}

貿易遵守とは、物品、サービス、技術の輸入、輸出、再輸出に関する法的制約を守ることである。ソフトウェア、暗号技術、クラウドサービス、開発ツール、技術文書なども対象になり得る。輸出管理、制裁対象国や企業、政府許可、輸入制限、関税などは専門性が高いため、実務では貿易遵守の専門家に確認する必要がある。

\paragraph*{サイバー犯罪}

\addcontentsline{toc}{subsubsection}{1.7.9 サイバー犯罪}

サイバー犯罪とは、コンピュータ、ソフトウェア、ネットワーク、組込みソフトウェアが関わる犯罪である。コンピュータやソフトウェアが犯罪の手段になる場合も、攻撃対象になる場合もある。詐欺、不正アクセス、盗聴、スパム、脅迫、嫌がらせ、個人データや営業秘密の窃取、他のシステムへの侵入や破壊などが含まれる。

ソフトウェア技術者には、サイバー犯罪の脅威を考慮し、ソフトウェアや利用者情報を偶発的または悪意あるアクセス、利用、変更、破壊、開示から守る責任がある。

\par\smallskip
\begin{center}
\centering
\IfFileExists{assets/generated_figures/ch14/fig-ch14-05.png}{\includegraphics[width=0.96\linewidth]{assets/generated_figures/ch14/fig-ch14-05.png}}{\fbox{\parbox[c][48mm][c]{0.9\linewidth}{\centering 画像生成待ち\\文書化の読者別情報を示す図}}}
\par\smallskip
\refstepcounter{figure}\label{fig:ch14-05}図 \thefigure: 文書化の読者別情報を示す図
\end{center}
図 \thefigure は、文書化の読者別情報を示す図を視覚的に整理したものである。
\par\smallskip

\subsubsection*{トレードオフ分析}

\addcontentsline{toc}{subsubsection}{1.9 トレードオフ分析}

トレードオフ分析とは、複数の解決案について、リスク、費用、便益を利害関係者と協力して評価し、目的に最も合う案を選ぶことである。ソフトウェアでは、性能、信頼性、安全性、セキュリティ、費用、納期、保守性、利用者体験などが競合しやすい。

プロジェクトが遅延したり予算超過したりした場合、どの要求を緩和または削除できるかを判断するためにトレードオフ分析が行われることがある。最初の手順は、設計目標を定め、その相対的な重要度を決めることである。目標の表現が曖昧だと、分析結果も曖昧になる。

トレードオフ分析は倫理的に行う必要がある。代替案を比較する基準と重みづけは、客観的かつ公平に選ぶべきである。利害の衝突がある場合は、あらかじめ開示しなければならない。

\par\smallskip
\begin{center}
\centering
\IfFileExists{assets/generated_figures/ch14/fig-ch14-06.png}{\includegraphics[width=0.96\linewidth]{assets/generated_figures/ch14/fig-ch14-06.png}}{\fbox{\parbox[c][48mm][c]{0.9\linewidth}{\centering 画像生成待ち\\トレードオフ分析の流れを示す図}}}
\par\smallskip
\refstepcounter{figure}\label{fig:ch14-06}図 \thefigure: トレードオフ分析の流れを示す図
\end{center}
図 \thefigure は、トレードオフ分析の流れを示す図を視覚的に整理したものである。
\par\smallskip
